% ============================================================================
% CHAPTER 12: PRODUCTION NLP AND DEPLOYMENT
% ============================================================================
% Covers tokenization algorithms, model serving infrastructure, quantization,
% distillation, speculative decoding, efficiency techniques, multilingual
% production challenges, and MLOps for NLP.
% ============================================================================

\chapter{Production NLP and Deployment}
\label{chap:production_nlp}

\begin{tldr}
Bridging the gap from research to production is where senior engineers differentiate
themselves.  This chapter covers the complete production NLP stack: tokenization algorithms
(BPE, WordPiece, Unigram), model serving infrastructure (KV-cache, continuous batching,
PagedAttention), efficiency techniques (quantization, distillation, speculative decoding),
and MLOps practices.  Every concept is grounded in concrete numbers---memory budgets,
latency targets, throughput benchmarks---because interviewers at L5+ expect you to reason
quantitatively about deployment tradeoffs, not just recite architectures.  Production
awareness in every answer signals staff-level thinking.
\end{tldr}

\bigskip

% ============================================================================
% SECTION 1: TOKENIZATION ALGORITHMS
% ============================================================================
\section{Tokenization Algorithms}
\label{sec:tokenization_algorithms}

Tokenization---splitting raw text into the discrete units a model consumes---is deceptively
one of the most consequential design decisions in any NLP system.  The choice of tokenizer
determines the vocabulary size, the sequence length (and therefore compute cost), the model's
ability to handle rare words and multiple languages, and even whether the model can perform
character-level tasks like spelling.

\subsection{Why Tokenization Matters}

\begin{enumerate}[label=\textbf{\arabic*.}]
    \item \textbf{Vocabulary size.}  Word-level tokenization produces vocabularies of
          hundreds of thousands of entries, most of which are rare.  Subword tokenization
          compresses this to 32K--100K tokens while retaining coverage of arbitrary text.
    \item \textbf{Sequence length.}  Every token costs compute ($O(n^2)$ in attention).
          A tokenizer that produces 30 tokens for the same sentence another produces 10
          tokens for is $9\times$ more expensive in attention FLOPs.
    \item \textbf{OOV handling.}  Word-level tokenizers assign a single \texttt{[UNK]}
          token to unseen words, discarding all information.  Subword tokenizers decompose
          unseen words into known pieces, preserving morphological signal.
    \item \textbf{Multilinguality.}  Tokenizer design directly determines how well a model
          handles non-English languages.  A tokenizer trained primarily on English will
          shatter CJK or Thai text into byte-level fragments, inflating sequence length.
\end{enumerate}

\begin{warningbox}
Interviewers use tokenization questions as a litmus test for production awareness.  A
candidate who cannot explain how BPE works, or who does not know that tokenizer choice
affects multilingual performance, reveals a gap between paper-reading knowledge and
system-building experience.
\end{warningbox}

\subsection{Byte Pair Encoding (BPE)}
\label{sec:bpe}

BPE (Sennrich et al., ACL 2016) is the most widely used subword tokenization algorithm
and the foundation for GPT-2, GPT-3, GPT-4, LLaMA, and Mistral.

\paragraph{Training algorithm.}
\begin{enumerate}
    \item \textbf{Initialize} the vocabulary with all individual characters (or bytes)
          present in the training corpus.
    \item \textbf{Count} the frequency of every adjacent symbol pair in the corpus.
    \item \textbf{Merge} the most frequent pair into a new symbol and add it to the
          vocabulary.
    \item \textbf{Repeat} steps 2--3 for a predefined number of merges (determining
          the final vocabulary size).
\end{enumerate}

\paragraph{Inference (encoding).}
Given a new word, apply the learned merge rules \emph{greedily} in the order they were
learned during training.  The word is first split into characters, then iteratively
merged according to the priority of each rule.

\paragraph{Handling OOV words.}
BPE never produces an \texttt{[UNK]} token (assuming byte-level BPE).  Any unseen word
is decomposed by splitting it into characters (or bytes) and applying the learned merges
in priority order; whatever subwords the merges produce are the output.  (Greedy
longest-match-first lookup is \emph{WordPiece} inference, not BPE.)  For example, a rare
word like ``unparalleledly'' might tokenize as \texttt{["un", "parallel", "ed", "ly"]}.

\begin{keyinsight}[BPE Is a Compression Algorithm]
BPE was originally a data compression algorithm (Gage, 1994).  Its application to NLP is
natural: the most frequent substrings get their own tokens, just as a compression
dictionary assigns short codes to frequent patterns.  This framing helps explain why
vocabulary size is a compression--compute tradeoff: larger vocabulary $=$ shorter
sequences (better compression) but larger embedding tables and softmax layers.
\end{keyinsight}

\subsection{WordPiece}
\label{sec:wordpiece}

WordPiece (Schuster and Nakajima, 2012; used in BERT, DistilBERT) is structurally
similar to BPE but differs in its merge criterion.

\paragraph{Key difference from BPE.}
BPE merges the pair with the highest \emph{frequency}.  WordPiece merges the pair that
maximizes the \emph{likelihood} of the training data under a unigram language model:
\[
    \text{score}(x, y) = \frac{\text{freq}(xy)}{\text{freq}(x) \cdot \text{freq}(y)}
\]
This prefers merges that create tokens which are more than the sum of their parts---pairs
that co-occur more than expected by chance.  In practice, the outputs of BPE and WordPiece
are similar, but WordPiece tends to produce slightly more linguistically meaningful tokens.

\paragraph{Special prefix.}
WordPiece marks continuation tokens with a \texttt{\#\#} prefix.  The word
``playing'' might tokenize as \texttt{["play", "\#\#ing"]}, explicitly indicating that
\texttt{\#\#ing} is not a word start.

\subsection{Unigram Language Model}
\label{sec:unigram}

The Unigram model (Kudo, EMNLP 2018) takes the opposite approach to BPE:

\paragraph{Training algorithm.}
\begin{enumerate}
    \item \textbf{Initialize} with a large seed vocabulary (e.g., all substrings up to
          a maximum length, or the top 1M most frequent substrings).
    \item \textbf{Estimate} token probabilities using EM: the E-step finds the most
          probable segmentation of each word; the M-step updates token probabilities
          from these segmentations.
    \item \textbf{Prune} tokens whose removal causes the least increase in corpus
          likelihood.  Remove a fixed percentage (e.g., 20\%) of the vocabulary.
    \item \textbf{Repeat} steps 2--3 until the target vocabulary size is reached.
\end{enumerate}

\paragraph{Inference (encoding).}
Unigram finds the \emph{most probable} segmentation using the Viterbi algorithm.  This
is a key advantage over BPE: Unigram has a principled probabilistic model over
segmentations, while BPE's greedy merge is a heuristic.

\paragraph{Used by:} T5, XLNet, ALBERT (all via SentencePiece).  Note that LLaMA also
uses SentencePiece, but with the \emph{BPE} algorithm, not Unigram.

\subsection{SentencePiece and Byte-Level BPE}

\paragraph{SentencePiece} (Kudo and Richardson, EMNLP 2018) is a tokenizer library
(not an algorithm) that implements both BPE and Unigram.  Its key feature is
\emph{language agnosticism}: it treats the input as a raw stream of Unicode text with no
pre-tokenization (no whitespace splitting, no language-specific rules), with optional
byte-fallback for characters outside the vocabulary.  This makes it
essential for multilingual models.  Whitespace is treated as a regular character
(represented as \texttt{\textunderscore}).

\paragraph{Byte-level BPE} (used in GPT-2+) operates on raw UTF-8 bytes rather than
Unicode characters.  Since there are only 256 possible byte values, the base vocabulary
is tiny and \emph{every possible text} can be represented---no \texttt{[UNK]} tokens,
ever.  This is critical for handling arbitrary scripts, code, and mixed-language text.

\subsection{Tokenization for Code}

Code has properties that natural language tokenizers handle poorly: significant whitespace
(Python indentation), long identifiers (\texttt{getUserAccountBalance}), operators, and
structured syntax.  Production approaches include:
\begin{itemize}
    \item \textbf{Whitespace normalization:} Replace runs of spaces with special tokens
          (e.g., Codex uses tokens for 1--24 spaces).
    \item \textbf{CamelCase/snake\_case splitting:} Some tokenizers split identifiers
          at case boundaries or underscores before applying BPE.
    \item \textbf{Larger vocabulary:} Code-focused models often use larger vocabularies
          (e.g., 100K+) to represent common code patterns as single tokens.
\end{itemize}

\subsection{Vocabulary Size Tradeoffs}

\begin{table}[H]
\centering
\small
\begin{tabularx}{\textwidth}{@{}l c c X@{}}
\toprule
\textbf{Vocab Size} & \textbf{Seq.\ Length} & \textbf{Embedding Params} & \textbf{Models} \\
\midrule
$\sim$30--32K & Longer   & 130M (at $d{=}4096$) & BERT (30,522), LLaMA 1/2 (32K) \\
$\sim$50K  & Moderate & 200M (at $d{=}4096$) & GPT-2/GPT-3 (50,257) \\
$\sim$100--128K & Shorter  & 400--520M (at $d{=}4096$) & GPT-4 ($\sim$100K), LLaMA 3 (128K) \\
$\sim$256K & Shortest & 1B (at $d{=}4096$) & Gemini, Gemma \\
\bottomrule
\end{tabularx}
\caption{Vocabulary size tradeoffs.  Larger vocabularies reduce sequence length (saving
compute in attention) but increase the embedding table size.  Modern models trend toward
larger vocabularies as the attention savings outweigh the embedding cost.}
\label{tab:vocab_size}
\end{table}

\subsection{Comparison of Tokenization Methods}

\begin{table}[H]
\centering
\small
\begin{tabularx}{\textwidth}{@{}l l l l X@{}}
\toprule
\textbf{Method} & \textbf{Merge Criterion} & \textbf{Direction} & \textbf{UNK?} & \textbf{Used By} \\
\midrule
BPE              & Most frequent pair        & Bottom-up & No (byte-level)  & GPT-2/3/4, LLaMA, Mistral \\
WordPiece        & Max likelihood gain       & Bottom-up & Possible         & BERT, DistilBERT \\
Unigram          & Min likelihood loss        & Top-down  & No (with bytes)  & T5, XLNet, ALBERT \\
SentencePiece    & Library (BPE or Unigram)  & Either    & No               & LLaMA, T5, mBART \\
Byte-level BPE   & Frequency (on bytes)       & Bottom-up & Never            & GPT-2+, all modern LLMs \\
\bottomrule
\end{tabularx}
\caption{Comparison of subword tokenization methods.}
\label{tab:tokenizer_comparison}
\end{table}

\subsection{TikZ: BPE Merge Process}

\begin{figure}[H]
\centering
\begin{tikzpicture}[
    node distance=0.6cm,
    tok/.style={rectangle, draw=deepblue, fill=lightblue, minimum width=0.7cm,
                minimum height=0.55cm, font=\footnotesize\ttfamily, rounded corners=2pt,
                thick, inner sep=3pt},
    mergetok/.style={rectangle, draw=deepgreen!80, fill=green!8, minimum width=0.7cm,
                minimum height=0.55cm, font=\footnotesize\ttfamily, rounded corners=2pt,
                thick, inner sep=3pt},
    arr/.style={-{Stealth[length=4pt]}, thick, deepblue!60},
    steplbl/.style={font=\small\sffamily\bfseries, text=deepblue},
    mergelbl/.style={font=\scriptsize\sffamily, text=deepgreen!80!black}
]

% Step 0: Characters
\node[steplbl] (s0) at (0, 0) {Init:};
\node[tok, right=0.3cm of s0] (c1) {l};
\node[tok, right=0.08cm of c1] (c2) {o};
\node[tok, right=0.08cm of c2] (c3) {w};
\node[tok, right=0.08cm of c3] (c4) {e};
\node[tok, right=0.08cm of c4] (c5) {s};
\node[tok, right=0.08cm of c5] (c6) {t};

% Step 1: Merge (e,s) -> es
\node[steplbl, below=0.7cm of s0] (s1) {Merge 1:};
\node[tok, right=0.3cm of s1] (m1a) {l};
\node[tok, right=0.08cm of m1a] (m1b) {o};
\node[tok, right=0.08cm of m1b] (m1c) {w};
\node[mergetok, right=0.08cm of m1c] (m1d) {es};
\node[tok, right=0.08cm of m1d] (m1e) {t};
\node[mergelbl, right=0.3cm of m1e] {(e, s) $\to$ es};

% Step 2: Merge (es,t) -> est
\node[steplbl, below=0.7cm of s1] (s2) {Merge 2:};
\node[tok, right=0.3cm of s2] (m2a) {l};
\node[tok, right=0.08cm of m2a] (m2b) {o};
\node[tok, right=0.08cm of m2b] (m2c) {w};
\node[mergetok, right=0.08cm of m2c] (m2d) {est};
\node[mergelbl, right=0.3cm of m2d] {(es, t) $\to$ est};

% Step 3: Merge (l,o) -> lo
\node[steplbl, below=0.7cm of s2] (s3) {Merge 3:};
\node[mergetok, right=0.3cm of s3] (m3a) {lo};
\node[tok, right=0.08cm of m3a] (m3b) {w};
\node[tok, right=0.08cm of m3b] (m3c) {est};
\node[mergelbl, right=0.3cm of m3c] {(l, o) $\to$ lo};

% Step 4: Merge (lo,w) -> low
\node[steplbl, below=0.7cm of s3] (s4) {Merge 4:};
\node[mergetok, right=0.3cm of s4] (m4a) {low};
\node[tok, right=0.08cm of m4a] (m4b) {est};
\node[mergelbl, right=0.3cm of m4b] {(lo, w) $\to$ low};

% Step 5: Merge (low,est) -> lowest
\node[steplbl, below=0.7cm of s4] (s5) {Merge 5:};
\node[mergetok, right=0.3cm of s5] (m5a) {lowest};
\node[mergelbl, right=0.3cm of m5a] {(low, est) $\to$ lowest};

% Arrows between steps
\draw[arr] ([yshift=-0.15cm]c3.south) -- ([yshift=0.15cm]m1c.north);
\draw[arr] ([yshift=-0.15cm]m1d.south) -- ([yshift=0.15cm]m2d.north);
\draw[arr] ([yshift=-0.15cm]m2a.south) -- ([yshift=0.15cm]m3a.north);
\draw[arr] ([yshift=-0.15cm]m3a.south) -- ([yshift=0.15cm]m4a.north);
\draw[arr] ([yshift=-0.15cm]m4a.south) -- ([yshift=0.15cm]m5a.north);

\end{tikzpicture}
\caption{BPE merge process for the word ``lowest.''  Starting from individual characters,
the algorithm iteratively merges the most frequent adjacent pair.  The merge order
becomes the encoding rule: at inference time, the same merges are applied greedily
to any input text.}
\label{fig:bpe_merge}
\end{figure}


% ============================================================================
% SECTION 2: MODEL SERVING INFRASTRUCTURE
% ============================================================================
\section{Model Serving Infrastructure}
\label{sec:model_serving}

Serving LLMs in production is fundamentally different from training them.  Training is
throughput-optimized (process as many tokens as possible per second); serving is
latency-optimized (minimize time-to-first-token and inter-token latency for each user).
This section covers the infrastructure primitives that make production LLM serving possible.

\subsection{KV-Cache: The Essential Optimization}
\label{sec:kv_cache}

In autoregressive generation, each new token attends to \emph{all} previous tokens.
Without caching, generating the $t$-th token requires recomputing the key and value
projections for all $t-1$ previous tokens---$O(t)$ redundant computation per step,
and $O(T^2)$ total for a sequence of length $T$.

The \textbf{KV-cache} stores the key and value tensors from all previous positions so
they are computed only once.

\paragraph{Two-phase generation.}
\begin{enumerate}
    \item \textbf{Prefill (prompt processing):}  Process all $n$ prompt tokens in
          parallel (like training).  Compute and store $\mK, \mV \in \R^{n \times d}$
          for every layer.  This is compute-bound.
    \item \textbf{Decode (token generation):}  Generate one token at a time.  For each
          new token, compute its query $\vq_t$, attend to the cached keys and values,
          then \emph{append} the new key and value to the cache.  This is
          memory-bandwidth-bound (reading the entire cache for a single token).
\end{enumerate}

\begin{keyinsight}[KV-Cache Is the Main Inference Bottleneck]
During decoding, the model reads the entire KV-cache from GPU memory for every single
token generated.  For a 70B parameter model with GQA ($H_{\text{kv}}{=}8$) at 128K
context, the KV-cache alone exceeds 40\,GB \emph{per request}---larger than the INT4-quantized
weights ($\sim$35\,GB), and across even a few concurrent long-context requests the total
cache outgrows the 140\,GB of FP16 weights.  This makes LLM inference
\textbf{memory-bandwidth-bound}, not compute-bound.  Understanding this distinction is
essential: most serving optimizations (GQA, quantization, PagedAttention) target
KV-cache memory, not raw compute.
\end{keyinsight}

\paragraph{KV-cache memory calculation.}
For a single request:
\[
    \text{KV memory} = 2 \times L \times H \times S \times d_h \times b
\]
where $L$ is the number of layers, $H$ is the number of KV heads, $S$ is the sequence
length, $d_h$ is the head dimension, $b$ is bytes per element (2 for FP16), and the
factor of 2 accounts for both keys and values.

\begin{mathresult}{KV-Cache Memory Budget}
\[
    M_{\text{KV}} = 2 \cdot L \cdot H_{\text{kv}} \cdot S \cdot d_h \cdot b
\]
\textbf{Concrete example} (LLaMA-2 7B, FP16):\\
$L = 32$, $H_{\text{kv}} = 32$ (MHA), $d_h = 128$, $b = 2$ bytes, $S = 2048$:
\[
    M_{\text{KV}} = 2 \times 32 \times 32 \times 2048 \times 128 \times 2
    = 1.07\;\text{GB per request}
\]
With GQA ($H_{\text{kv}} = 8$ instead of 32): $M_{\text{KV}} = 268\;\text{MB}$---a $4\times$ reduction.\\
At 128K context: $M_{\text{KV}} = 67\;\text{GB}$ (MHA) or $16.8\;\text{GB}$ (GQA).
\end{mathresult}

\subsection{TikZ: KV-Cache During Autoregressive Decoding}

\begin{figure}[H]
\centering
\begin{tikzpicture}[
    node distance=0.5cm and 0.4cm,
    cachebox/.style={rectangle, draw=deepblue, fill=lightblue, minimum width=0.6cm,
                     minimum height=0.6cm, font=\scriptsize\sffamily, thick},
    newbox/.style={rectangle, draw=deepgreen!80, fill=green!12, minimum width=0.6cm,
                   minimum height=0.6cm, font=\scriptsize\sffamily, thick},
    emptybox/.style={rectangle, draw=gray!40, fill=gray!5, minimum width=0.6cm,
                     minimum height=0.6cm, font=\scriptsize\sffamily, dashed},
    lbl/.style={font=\small\sffamily\bfseries, text=deepblue},
    sublbl/.style={align=center, font=\scriptsize\sffamily, text=gray!70!black},
    arr/.style={-{Stealth[length=4pt]}, thick, deepblue!60},
    brace/.style={decorate, decoration={brace, amplitude=5pt, raise=2pt}, thick, deepblue!60}
]

% Labels
\node[lbl] at (-1.5, 2.5) {Step};
\node[lbl] at (2.0, 2.5) {KV-Cache Contents};
\node[lbl] at (6.5, 2.5) {New Token};

% Step 1: Prefill
\node[sublbl] at (-1.5, 1.5) {Prefill};
\node[cachebox] (p1) at (0.5, 1.5) {$k_1$};
\node[cachebox, right=0.08cm of p1] (p2) {$k_2$};
\node[cachebox, right=0.08cm of p2] (p3) {$k_3$};
\node[cachebox, right=0.08cm of p3] (p4) {$k_4$};
\node[emptybox, right=0.08cm of p4] (pe) {};
\node[emptybox, right=0.08cm of pe] (pf) {};
\node[sublbl, right=0.3cm of pf] {(compute-bound)};

% Step 2: Decode t=5
\node[sublbl] at (-1.5, 0.5) {Decode $t{=}5$};
\node[cachebox] (d1) at (0.5, 0.5) {$k_1$};
\node[cachebox, right=0.08cm of d1] (d2) {$k_2$};
\node[cachebox, right=0.08cm of d2] (d3) {$k_3$};
\node[cachebox, right=0.08cm of d3] (d4) {$k_4$};
\node[newbox, right=0.08cm of d4] (d5) {$k_5$};
\node[emptybox, right=0.08cm of d5] (de) {};
\node[newbox] at (6.5, 0.5) {$t_5$};

% Step 3: Decode t=6
\node[sublbl] at (-1.5, -0.5) {Decode $t{=}6$};
\node[cachebox] (e1) at (0.5, -0.5) {$k_1$};
\node[cachebox, right=0.08cm of e1] (e2) {$k_2$};
\node[cachebox, right=0.08cm of e2] (e3) {$k_3$};
\node[cachebox, right=0.08cm of e3] (e4) {$k_4$};
\node[cachebox, right=0.08cm of e4] (e5) {$k_5$};
\node[newbox, right=0.08cm of e5] (e6) {$k_6$};
\node[newbox] at (6.5, -0.5) {$t_6$};

% Arrows showing cache growth
\draw[arr] (d5.south) -- ([yshift=0.15cm]e5.north);
\draw[arr, deepgreen!70] (6.5, 0.2) -- (6.5, -0.2);

% Annotation
\node[sublbl, align=left] at (3.0, -1.3) {Cache grows by one K/V pair per decode step.\\
Each step reads the \emph{entire} cache (memory-bandwidth-bound).};

\end{tikzpicture}
\caption{KV-cache during autoregressive generation.  The prefill phase computes keys
and values for all prompt tokens in parallel.  Each subsequent decode step appends one
new K/V pair and reads the full cache to compute attention over all previous tokens.}
\label{fig:kv_cache}
\end{figure}

\subsection{Continuous Batching}
\label{sec:continuous_batching}

Traditional \textbf{static batching} pads all sequences in a batch to the same length
and waits for the longest sequence to finish before starting a new batch.  If one
request generates 10 tokens and another generates 500, the short request's GPU slot is
wasted for 490 steps.

\textbf{Continuous batching} (also called \emph{iteration-level scheduling}) solves this:
\begin{itemize}
    \item At \emph{every} decode step, check for completed sequences and newly arrived
          requests.
    \item Remove completed sequences immediately, freeing their KV-cache memory.
    \item Insert new requests into the batch without waiting for others to finish.
    \item Result: GPU utilization approaches 100\%, throughput increases 2--8$\times$
          over static batching in practice.
\end{itemize}

This is the default scheduling strategy in production serving frameworks (vLLM,
TensorRT-LLM, TGI).

\subsection{PagedAttention (vLLM)}
\label{sec:paged_attention}

The KV-cache is typically allocated as a contiguous block per request.  Since the final
sequence length is unknown in advance, systems over-allocate, wasting 60--80\% of
KV-cache memory on average.

\textbf{PagedAttention} (Kwon et al., SOSP 2023) applies the operating system concept
of \emph{virtual memory} to KV-cache management:
\begin{itemize}
    \item KV-cache is divided into fixed-size \textbf{blocks} (e.g., 16 tokens per block).
    \item Blocks are allocated \emph{on demand} as the sequence grows, from a shared
          pool---no contiguous allocation required.
    \item A \textbf{block table} (analogous to a page table) maps each sequence's
          logical token positions to physical memory blocks.
    \item \textbf{Result:} Near-zero memory waste.  Memory utilization improves from
          $\sim$20\% to $\sim$95\%, enabling 2--4$\times$ higher batch sizes.
\end{itemize}

An additional benefit: sequences that share a common prefix (e.g., the same system
prompt) can share the same physical blocks via \textbf{copy-on-write}, dramatically
reducing memory for multi-turn conversations and batched requests with identical
system prompts.

\subsection{Parallelism Strategies}
\label{sec:parallelism}

When a model does not fit on a single GPU, or when throughput requirements demand
multiple GPUs, there are two primary parallelism strategies:

\begin{table}[H]
\centering
\small
\begin{tabularx}{\textwidth}{@{}l X X@{}}
\toprule
\textbf{Strategy} & \textbf{Tensor Parallelism (TP)} & \textbf{Pipeline Parallelism (PP)} \\
\midrule
What is split & Individual layers split across GPUs & Different layers on different GPUs \\
Communication & All-reduce after every layer & Point-to-point between stages \\
Latency       & Low (all GPUs process each token) & Higher (pipeline bubbles) \\
Throughput    & Limited by communication bandwidth & High with micro-batching \\
Typical use   & Within a single node (fast NVLink) & Across nodes (slower network) \\
\bottomrule
\end{tabularx}
\caption{Tensor parallelism vs.\ pipeline parallelism for LLM serving.}
\label{tab:parallelism}
\end{table}

In practice, production deployments often combine both: TP within a node (e.g., 8 GPUs
connected by NVLink) and PP across nodes.

\subsection{Prefix Caching}
\label{sec:prefix_caching}

Many production workloads share a common prefix: the same system prompt is prepended to
every user request.  \textbf{Prefix caching} precomputes the KV-cache for this shared
prefix once and reuses it for all requests:
\begin{itemize}
    \item For a 2,000-token system prompt on LLaMA-2 7B, prefix caching saves
          $\sim$1\,GB of KV-cache memory per request and eliminates
          $\sim$100\,ms of prefill latency.
    \item This is implemented in vLLM (automatic prefix caching), TGI, and
          TensorRT-LLM.
    \item Combined with PagedAttention's copy-on-write, prefix caching enables
          serving thousands of concurrent requests with a shared system prompt.
\end{itemize}


% ============================================================================
% SECTION 3: QUANTIZATION
% ============================================================================
\section{Quantization}
\label{sec:quantization}

Quantization reduces the numerical precision of model weights and/or activations,
trading a small amount of quality for significant improvements in memory, speed, and cost.
For LLMs, quantization is often the difference between a model fitting on available
hardware or not.

\subsection{Fundamentals}

\begin{mathresult}{Uniform Quantization}
A floating-point value $x$ is quantized to an integer representation:
\[
    q = \text{round}\!\left(\frac{x}{s} + z\right), \qquad
    \hat{x} = s \cdot (q - z)
\]
where $s$ is the \textbf{scale factor} (maps the floating-point range to integer range),
$z$ is the \textbf{zero point} (ensures zero maps exactly), and $\hat{x}$ is the
dequantized approximation.  The quantization error is $|x - \hat{x}|$.
\end{mathresult}

\paragraph{Key precision levels.}
\begin{itemize}
    \item \textbf{FP16 / BF16 (16-bit):}  Standard training and inference precision.
          BF16 preferred for training (larger dynamic range, same as FP32 exponent).
          No quality loss versus FP32 for most models.
    \item \textbf{INT8 (8-bit):}  $2\times$ memory reduction, $\sim$$2\times$ speedup.
          Per-channel weight quantization with per-token activation quantization yields
          minimal quality loss ($<$1\% on most benchmarks).
    \item \textbf{INT4 (4-bit):}  $4\times$ memory reduction.  Requires careful
          quantization methods (GPTQ, AWQ) to preserve quality.  Typical quality loss:
          1--3\% on standard benchmarks.
    \item \textbf{FP8:}  Hardware-supported on NVIDIA H100 and later.  Combines the
          simplicity of floating-point with 8-bit efficiency.  Increasingly the default
          for inference.
\end{itemize}

\subsection{Post-Training Quantization (PTQ)}

PTQ quantizes a pre-trained model without retraining.  The challenge is minimizing the
quantization error, which is uneven across weights---a few ``outlier'' weights with large
magnitude contribute disproportionately to the output.

\paragraph{GPTQ} (Frantar et al., ICLR 2023) performs \textbf{layer-wise quantization}
using second-order information:
\begin{itemize}
    \item Quantize weights one column at a time.
    \item For each quantized column, adjust the \emph{remaining} unquantized columns to
          compensate for the error, using the inverse Hessian of the layer's output.
    \item This ``error compensation'' preserves the layer's output much better than
          naive rounding.
    \item \textbf{Result:} 4-bit weight quantization with near-FP16 quality.
          A 70B model fits on a single 48\,GB GPU (vs.\ 140\,GB in FP16).
\end{itemize}

\paragraph{AWQ} (Lin et al., MLSys 2024) takes a different approach---\textbf{activation-aware}
quantization:
\begin{itemize}
    \item Observation: not all weights are equally important.  Weights corresponding
          to large-magnitude activations matter more.
    \item Scale up the important weight channels before quantization (equivalent to
          finding the optimal per-channel scale that minimizes quantization error
          weighted by activation magnitude).
    \item \textbf{Advantage over GPTQ:} Simpler, faster to apply, and often produces
          better quality at 4-bit because it directly accounts for which weights
          the model relies on most.
\end{itemize}

\subsection{Quantization-Aware Training (QAT)}

QAT simulates quantization during training by inserting ``fake quantization'' operations
in the forward pass: weights are quantized and immediately dequantized, so the forward
pass sees quantized values but gradients still flow through (using the straight-through
estimator).  QAT consistently outperforms PTQ but requires full training infrastructure,
making it expensive for large models.

\subsection{Quality--Speed--Memory Tradeoffs}

\begin{table}[H]
\centering
\small
\begin{tabularx}{\textwidth}{@{}l c c c c X@{}}
\toprule
\textbf{Precision} & \textbf{Mem.\ Saving} & \textbf{Speed} & \textbf{Quality} & \textbf{Method} & \textbf{Notes} \\
\midrule
FP16/BF16 & $1\times$ (baseline) & $1\times$    & Baseline    & None  & Standard \\
INT8      & $2\times$            & $\sim$$2\times$ & $<$1\% drop & LLM.int8()  & Safe default \\
FP8       & $2\times$            & $\sim$$2\times$ & $<$0.5\% drop & Native H100 & Best 8-bit \\
INT4 (GPTQ) & $4\times$         & $\sim$$3\times$ & 1--3\% drop & GPTQ  & Hessian-aware \\
INT4 (AWQ) & $4\times$          & $\sim$$3\times$ & 1--2\% drop & AWQ   & Activation-aware \\
INT3/INT2 & $5$--$8\times$      & $\sim$$4\times$ & 5--15\% drop & Experimental & Research only \\
\bottomrule
\end{tabularx}
\caption{Quantization precision levels with typical tradeoffs.  Numbers are approximate
and vary significantly by model size and task.}
\label{tab:quantization_tradeoffs}
\end{table}

\begin{warningbox}
Quantization quality loss varies dramatically across models, tasks, and evaluation
benchmarks.  A 4-bit quantized model might show only 1\% degradation on MMLU but 10\%
on a math reasoning benchmark.  \emph{Always evaluate quantized models on your specific
task before deploying.}  The ``1--3\% quality loss'' numbers widely cited are averages
over easy benchmarks and may not reflect your use case.
\end{warningbox}


% ============================================================================
% SECTION 4: KNOWLEDGE DISTILLATION
% ============================================================================
\section{Knowledge Distillation}
\label{sec:distillation}

Knowledge distillation (Hinton et al., 2015) trains a small \textbf{student} model to
mimic the behavior of a large \textbf{teacher} model.  The key insight is that the
teacher's \emph{soft probability distribution} over classes contains far more information
than the hard labels alone.

\paragraph{Why soft labels are more informative.}
When a teacher classifies an image of a car, its softmax output might be
$[0.85, 0.10, 0.03, 0.02]$ over $[\text{car}, \text{truck}, \text{bus},
\text{bicycle}]$.  The hard label says only ``car,'' but the soft distribution reveals
that this image is somewhat truck-like and slightly bus-like---\emph{inter-class
similarity} information that hard labels destroy.

\begin{mathresult}{Distillation Loss}
The student is trained with a weighted combination of two losses:
\[
    \loss = \alpha \cdot \loss_{\text{soft}} + (1 - \alpha) \cdot \loss_{\text{hard}}
\]
where:
\begin{align*}
    \loss_{\text{soft}} &= T^2 \cdot \KL\!\Big(\softmax\!\big(\vz_t / T\big) \;\big\|\; \softmax\!\big(\vz_s / T\big)\Big) \\
    \loss_{\text{hard}} &= \text{CrossEntropy}\!\big(\vy,\; \softmax(\vz_s)\big)
\end{align*}
$\vz_t$ and $\vz_s$ are the teacher and student logits, $T$ is the \textbf{temperature}
(typically 2--20), and $\alpha$ balances the two losses (typically 0.5--0.9).  The $T^2$
factor corrects for the gradient magnitude reduction when dividing logits by $T$.
\end{mathresult}

\paragraph{Temperature intuition.}
At $T = 1$ (standard softmax), the teacher's distribution is peaky---the correct class
dominates.  As $T$ increases, the distribution softens, revealing more of the relative
structure among classes.  Higher $T$ transfers more ``dark knowledge'' but makes the
signal noisier.

\subsection{Distillation in NLP}

\paragraph{DistilBERT} (Sanh et al., 2019): 6-layer student distilled from 12-layer
BERT-base.  \textbf{Result:} 60\% the size, 60\% faster inference, retaining 97\% of
BERT's performance on GLUE.  The distillation used a combination of language modeling
loss, distillation loss, and cosine embedding loss (aligning hidden states).

\paragraph{Modern LLM distillation.}
Distilling from large LLMs (e.g., GPT-4, Claude) into smaller open models has become
a standard practice:
\begin{itemize}
    \item \textbf{Output distillation:} Generate responses from the teacher, train the
          student on these (input, response) pairs via standard SFT.  This is the
          simplest form and does not require access to teacher logits.
    \item \textbf{Logit distillation:} When teacher logits are available, use the
          full KL-divergence loss.  Produces better students but requires access to
          the teacher's internals.
    \item \textbf{On-policy distillation:} The student generates outputs, the teacher
          scores or corrects them, and the student trains on the corrections.
          Avoids the distribution mismatch of pure offline distillation.
\end{itemize}

\paragraph{Phi models} (Microsoft, 2023--2024) demonstrated that training on
``textbook quality'' synthetic data generated by a larger model can produce small models
(1.3B--14B parameters) with outsized performance---a form of implicit distillation
through data curation.

\subsection{Distillation vs.\ Quantization vs.\ Pruning}

\begin{table}[H]
\centering
\small
\begin{tabularx}{\textwidth}{@{}l X X X@{}}
\toprule
& \textbf{Distillation} & \textbf{Quantization} & \textbf{Pruning} \\
\midrule
\textbf{What changes}    & Architecture (fewer layers) & Numerical precision & Sparsity (remove weights) \\
\textbf{Training needed}  & Yes (full student training) & No (PTQ) or minimal (QAT) & Minimal to moderate \\
\textbf{Typical speedup}  & $2$--$6\times$               & $2$--$4\times$              & $1.5$--$3\times$ \\
\textbf{Quality}          & 90--97\% of teacher          & 95--99\% of original        & 90--98\% of original \\
\textbf{Combinable?}      & Yes (distill then quantize)  & Yes (quantize anything)     & Yes (prune then quantize) \\
\bottomrule
\end{tabularx}
\caption{Comparison of model compression techniques.  In practice, these are often
combined: distill a large model into a smaller one, then quantize the student to INT4.}
\label{tab:compression_comparison}
\end{table}


% ============================================================================
% SECTION 5: SPECULATIVE DECODING
% ============================================================================
\section{Speculative Decoding}
\label{sec:speculative_decoding}

Autoregressive generation is inherently sequential: each token depends on the previous
one.  For a 70B parameter model, generating a single token takes $\sim$70\,ms of pure
memory traffic on A100-class hardware (reading 140\,GB of weights at $\sim$2\,TB/s
aggregate memory bandwidth across the GPUs holding the shards).  With speculative
decoding, we can generate 2--3$\times$ more tokens per second without any change in
output quality.

\subsection{The Algorithm}

\begin{enumerate}
    \item \textbf{Draft:} A small, fast \emph{draft model} (e.g., 1B parameters)
          autoregressively generates $k$ candidate tokens ($k$ is typically 4--8).
          This is cheap: $\sim$2\,ms per token.
    \item \textbf{Verify:} The large \emph{target model} processes all $k$ draft tokens
          \emph{in a single forward pass} (parallel, like prefill).  This produces
          the target model's probability distribution at each position.
    \item \textbf{Accept/reject:} Compare the draft model's distribution to the target
          model's distribution at each position, left to right.  Accept tokens as long
          as the draft distribution ``matches'' the target distribution; reject from the
          first mismatch onward.  Sample a correction token from a modified distribution
          at the rejection point.
\end{enumerate}

\begin{keyinsight}[Speculative Decoding Is ``Free'' Quality]
The accept/reject step is designed so that the \emph{final output distribution is
mathematically identical} to sampling from the target model alone.  There is zero quality
loss---the speedup is obtained purely by replacing sequential target-model calls with a
combination of cheap sequential draft-model calls and a single parallel target-model call.
This is provable, not just empirical.
\end{keyinsight}

\paragraph{Why it works: verification is cheaper than generation.}
Generating $k$ tokens sequentially with the target model requires $k$ forward passes
(each reading the full model weights from memory).  Verifying $k$ tokens requires only
\emph{one} forward pass (processing all $k$ tokens in parallel), with compute cost
roughly equal to generating a single token.  The speedup comes from this asymmetry
between sequential generation and parallel verification.

\paragraph{Acceptance rate.}
If the draft model closely approximates the target model, most tokens are accepted.
In practice, acceptance rates of 70--90\% are typical for well-matched draft-target
pairs.  The effective speedup is approximately:
\[
    \text{Speedup} \approx \frac{\text{avg.\ accepted tokens} + 1}{1 + k \cdot (c_{\text{draft}} / c_{\text{target}})}
\]
where $c_{\text{draft}} / c_{\text{target}} \ll 1$.  Worked through: with $k = 5$, a
per-token acceptance rate of $0.8$ (expected accepted tokens
$\sum_{i=1}^{5} 0.8^i \approx 2.7$), and a cost ratio
$c_{\text{draft}}/c_{\text{target}} \approx 0.03$, the formula gives
$(2.7 + 1)/(1 + 5 \times 0.03) \approx 3.2\times$; scheduling overhead and imperfect
batching bring realized speedups down to the commonly quoted $2$--$3\times$.

\paragraph{Variants.}
\begin{itemize}
    \item \textbf{Self-speculative decoding:} Use the target model itself as the draft
          by skipping layers (early exit).  No separate draft model needed.
    \item \textbf{Medusa:} Add multiple prediction heads to the target model that predict
          future tokens in parallel.  Train the extra heads; no separate model needed.
    \item \textbf{EAGLE:} Train a lightweight feature-level draft model on top of the
          target model's hidden states.  Higher acceptance rates than token-level drafting.
\end{itemize}


% ============================================================================
% SECTION 6: ADDITIONAL EFFICIENCY TECHNIQUES
% ============================================================================
\section{Additional Efficiency Techniques}
\label{sec:efficiency}

\subsection{Pruning}

Pruning removes unnecessary weights, heads, or layers from a trained model.

\paragraph{Unstructured pruning.}
Set individual weights to zero based on magnitude.  Can achieve 90\%+ sparsity with
minimal quality loss, but \emph{requires hardware that accelerates sparse operations}
(sparse tensor cores, N:M sparsity support on Ampere+).  Without hardware support,
unstructured sparsity provides no speedup---the zero weights are still stored and
multiplied.

\paragraph{Structured pruning.}
Remove entire attention heads, FFN neurons, or even full layers.  The resulting model
is a standard dense model with fewer parameters---no special hardware needed.  Typical
findings: 30--50\% of attention heads can be removed with $<$2\% quality loss
(Michel et al., NeurIPS 2019).

\subsection{Flash Attention}

Flash Attention (Dao et al., NeurIPS 2022; covered in detail in Chapter~\ref{chap:transformers})
is an \textbf{IO-aware} exact attention implementation that avoids materializing the
full $n \times n$ attention matrix in HBM:
\begin{itemize}
    \item Tiles the Q, K, V matrices into blocks that fit in SRAM (on-chip cache).
    \item Computes attention block-by-block using online softmax (numerically stable
          incremental computation).
    \item \textbf{Result:} $2$--$4\times$ speedup and $O(n)$ memory instead of $O(n^2)$.
    \item Flash Attention 2 adds further optimizations (better work partitioning,
          reduced non-matmul FLOPs).
    \item This is now the default attention implementation in PyTorch 2.0+
          (\texttt{scaled\_dot\_product\_attention}).
\end{itemize}

\subsection{Mixed Precision Training and Inference}

\begin{itemize}
    \item \textbf{FP16:}  Half the memory of FP32, hardware-accelerated on all modern GPUs.
          Requires loss scaling to handle the reduced dynamic range.
    \item \textbf{BF16:}  Same exponent range as FP32 (8 bits) with reduced mantissa
          (7 bits vs.\ 23).  No loss scaling needed.  The standard for LLM training.
    \item \textbf{FP8:}  Supported on H100+.  Two formats: E4M3 (4 exponent, 3 mantissa---better
          for weights) and E5M2 (5 exponent, 2 mantissa---better for gradients/activations).
\end{itemize}

\subsection{Gradient Checkpointing}

During training, all intermediate activations are stored for the backward pass.  For a
70B model with 2K context, this can exceed 100\,GB.  \textbf{Gradient checkpointing}
(also called activation recomputation) stores only a subset of activations (e.g., one
per layer) and recomputes the rest during the backward pass.  This trades $\sim$$33\%$
extra compute for a $\sim$$5\times$ reduction in activation memory.

\subsection{Model Compilation}

\begin{itemize}
    \item \textbf{\texttt{torch.compile}:} PyTorch 2.0's compiler.  Traces the
          computation graph, applies operator fusion, and generates optimized kernels.
          Typical speedup: 20--40\% for inference, with no code changes.
    \item \textbf{TensorRT / TensorRT-LLM:} NVIDIA's inference compiler.  Fuses operations,
          selects optimal kernels, applies quantization.  Aggressive optimization for
          NVIDIA hardware, 2--5$\times$ speedup over naive PyTorch.
    \item \textbf{ONNX Runtime:} Cross-platform inference optimization.  Useful for
          deploying on diverse hardware (CPU, mobile, edge).
\end{itemize}


% ============================================================================
% SECTION 7: MULTILINGUAL PRODUCTION CHALLENGES
% ============================================================================
\section{Multilingual Production Challenges}
\label{sec:multilingual_production}

Deploying NLP systems globally introduces challenges that monolingual (English-focused)
development obscures.

\subsection{The Tokenizer Tax}
\label{sec:tokenizer_tax}

Tokenizers trained primarily on English text produce dramatically more tokens for
non-English and non-Latin-script languages.  This ``tokenizer tax'' has cascading
effects:

\begin{itemize}
    \item \textbf{Effective context window is shorter.}  If Thai text tokenizes at
          $3\times$ the rate of English, a 128K-token context window holds
          $\sim$43K tokens worth of Thai content.
    \item \textbf{Inference is more expensive.}  More tokens $=$ more compute
          ($O(n^2)$ attention) and more KV-cache memory.
    \item \textbf{API costs increase.}  Token-based pricing penalizes non-English users.
    \item \textbf{Generation is slower.}  More tokens to generate for the same semantic
          content.
\end{itemize}

\paragraph{Concrete example.}
The sentence ``Artificial intelligence is transforming the world'' is approximately:
\begin{itemize}
    \item \textbf{English:} 7 tokens (GPT-4 tokenizer)
    \item \textbf{Chinese:} 11 tokens ($1.6\times$)
    \item \textbf{Japanese:} 14 tokens ($2.0\times$)
    \item \textbf{Thai:} 25 tokens ($3.6\times$)
    \item \textbf{Hindi:} 30 tokens ($4.3\times$)
\end{itemize}

\paragraph{Mitigation strategies.}
\begin{itemize}
    \item Train tokenizers on balanced multilingual corpora (equal representation of
          languages).
    \item Increase vocabulary size (100K+ tokens) to include more non-English subwords.
    \item Use byte-level BPE to ensure universal coverage.
    \item Evaluate fertility (tokens per word) across target languages as a tokenizer
          quality metric.
\end{itemize}

\subsection{Cross-Lingual Transfer}

Multilingual models (mBERT, XLM-R, NLLB-200) enable training on high-resource languages
and deploying on low-resource languages.  This works because shared subword vocabularies
and shared model parameters force the model to learn \emph{language-agnostic}
representations.

\paragraph{When cross-lingual transfer fails.}
\begin{itemize}
    \item Script mismatch: Latin-script training data transfers poorly to Devanagari or
          Thai script without explicit cross-script signal.
    \item Morphological mismatch: English (analytic) to Turkish (agglutinative)
          transfer loses morphological structure.
    \item Domain mismatch: legal or medical terminology is language-specific and does
          not transfer.
\end{itemize}

\subsection{Code-Switching}

Real-world text frequently mixes languages: ``Let's meet at the caf\'{e}, va bene?''
Code-switched text breaks assumptions of both language detection and monolingual models.
Production approaches include multilingual models (which handle mixed-language input
natively), language-aware tokenization, and training on code-switched data.


% ============================================================================
% SECTION 8: MLOPS FOR NLP
% ============================================================================
\section{MLOps for NLP}
\label{sec:mlops_nlp}

Deploying a model is the beginning, not the end.  Production NLP systems require
continuous monitoring, evaluation, and improvement.

\subsection{Monitoring}

\paragraph{Latency metrics.}
\begin{itemize}
    \item \textbf{Time to first token (TTFT):}  From request arrival to first generated
          token.  Dominated by prefill time.  Target: $<$500\,ms for interactive
          applications, $<$100\,ms for search.
    \item \textbf{Inter-token latency (ITL):}  Time between consecutive generated tokens.
          Determines perceived streaming speed.  Target: $<$50\,ms (20+ tokens/sec)
          for smooth reading experience.
    \item \textbf{End-to-end latency:}  Total time from request to complete response.
          P50 is the median; P95 and P99 capture tail latency.  SLAs are typically
          set on P99.
    \item \textbf{Throughput:}  Tokens per second (total across all requests) or
          requests per second.  The primary capacity planning metric.
\end{itemize}

\paragraph{Quality metrics.}
\begin{itemize}
    \item \textbf{Online evaluation:} User feedback (thumbs up/down), click-through
          rates, task completion rates, session length.
    \item \textbf{Automated evaluation:} Run shadow traffic through evaluation pipelines
          (LLM-as-judge, factuality checkers, toxicity classifiers).
    \item \textbf{Regression detection:} Maintain a golden test set; run every model
          update against it and block deployment on quality regression.
\end{itemize}

\subsection{Data Drift Detection}

Input distributions shift over time: new topics emerge, user demographics change,
adversarial inputs evolve.  Monitor for:
\begin{itemize}
    \item \textbf{Input length distribution:} Sudden changes may indicate bot traffic
          or new use patterns.
    \item \textbf{Token distribution:} Track KL divergence between current input token
          frequencies and the training distribution.
    \item \textbf{Embedding drift:} Monitor the centroid and spread of input embeddings;
          significant movement indicates distribution shift.
    \item \textbf{Confidence distribution:} A shift toward lower (or uniformly high)
          confidence scores signals that the model is encountering unfamiliar inputs.
\end{itemize}

\subsection{Cost Optimization}

\begin{itemize}
    \item \textbf{Model routing:}  Use a small classifier to route easy queries to a
          small model (7B, INT4) and hard queries to a large model (70B, FP16).  This
          can reduce average cost by 5--10$\times$ while maintaining quality on
          difficult queries.
    \item \textbf{Semantic caching:}  Cache responses for semantically similar queries.
          If a new query's embedding is within a threshold distance of a cached query,
          return the cached response.  Effective for FAQ-like workloads.
    \item \textbf{Prompt compression:}  Summarize or truncate long system prompts
          to reduce prefill cost.
    \item \textbf{Batching:}  Continuous batching maximizes GPU utilization; larger
          batch sizes amortize the fixed cost of reading model weights.
\end{itemize}

\subsection{Model Versioning and A/B Testing}

\begin{itemize}
    \item \textbf{Model registry:}  Version every model artifact (weights, tokenizer,
          config, evaluation results) with immutable identifiers.
    \item \textbf{A/B testing:}  Split traffic between model versions; measure
          quality and latency with statistical significance before full rollout.
    \item \textbf{Canary deployment:}  Route 1--5\% of traffic to the new model;
          monitor error rates and latency; gradually increase if metrics are stable.
    \item \textbf{Rollback:}  Maintain the ability to instantly revert to the previous
          model version if quality degrades.
\end{itemize}

\subsection{Feedback Loops}

\begin{enumerate}
    \item \textbf{Implicit feedback:}  User behavior signals (regenerate button clicks,
          copy-paste rates, session abandonment).
    \item \textbf{Explicit feedback:}  Thumbs up/down, star ratings, corrections.
    \item \textbf{Active learning:}  Route low-confidence examples to human annotators;
          retrain on the expanded dataset.
    \item \textbf{RLHF/DPO loops:}  Collect preference data from production; retrain
          reward models or run DPO on the preference data.
\end{enumerate}


% ============================================================================
% SECTION 9: INTERVIEW QUESTIONS
% ============================================================================
\section{Interview Questions}
\label{sec:production_interview_questions}

% --- Question 1: Serving at Scale ---
\begin{interviewq}
\textbf{Q: How would you serve an LLM at 10K QPS with less than 500ms time-to-first-token?}

\badgeLsix{L6} \quad \badgetype{System Design} \quad \badgefreq{\freqhigh}

\stronganswer
This requires reasoning across the full serving stack.  10K QPS with 500ms TTFT is
extremely demanding---let me state assumptions first, then work the token budget
honestly.  \textbf{Assumptions:} 7B-class model with GQA ($H_{\text{kv}}{=}8$),
1K-token prompts, $\sim$200-token responses, H100-class GPUs ($\sim$3.35\,TB/s HBM;
$\sim$990 dense BF16 TFLOPS of which $\sim$50\% is realistically sustained, i.e.\
$\sim$500 effective TFLOPS; FP8 roughly doubles that).

\textbf{Model selection and optimization:}  Start with the smallest model that meets
quality requirements.  A 7B model with INT4 weight quantization (AWQ or GPTQ) fits on
a single GPU with ample room for KV-cache.  TTFT is dominated by prefill, which is
\emph{compute}-bound: a 1K-token prompt costs $2N \approx 14$\,GFLOPs per token
$\times$ 1K tokens $= 14$\,TFLOPs, i.e.\ $\sim$30\,ms at 500 effective TFLOPS
($\sim$15\,ms with FP8)---well inside the 500ms budget \emph{per request}.  The hard
part is aggregate throughput, not single-request latency.  If a 70B model is required,
use GQA + quantization + tensor parallelism across a node, and every count below
scales up by roughly $10\times$.

\textbf{Batching and the token budget:}  Continuous batching is essential, but the GPU
count follows from the \emph{token} budget, not the request count.  10K QPS $\times$
1K prompt tokens $=$ 10M prefill tokens/s.  One H100 sustains roughly
$500\,\text{TFLOPS} / 14\,\text{GFLOPs per token} \approx 35\text{K}$ prefill tokens/s
in BF16 ($\sim$70K with FP8)---only $\sim$35--70 QPS of 1K-token prefill per GPU.
Prefill alone therefore needs $\sim$140 (FP8) to $\sim$280 (BF16) GPUs.  Decode adds
10K QPS $\times$ 200 output tokens $=$ 2M tokens/s; at high batch, decode is also
compute-bound at the same per-GPU token rate, adding another $\sim$30--60 GPUs.
Total: on the order of \textbf{200--350 H100s} ($\sim$25--45 8-GPU nodes), plus
headroom for P99 and traffic spikes---an order of magnitude more than a naive
request-count estimate suggests.

\textbf{Infrastructure:}
\begin{itemize}
    \item Use vLLM or TensorRT-LLM as the serving engine (PagedAttention, continuous
          batching, prefix caching).
    \item Prefix caching for shared system prompts eliminates redundant prefill.
    \item Load balance across GPU replicas with least-connections routing.
    \item Horizontal scaling: add replicas linearly to increase throughput.
\end{itemize}

\textbf{Latency optimization:}  Prefix caching reduces TTFT for repeated prefixes.
Speculative decoding reduces inter-token latency (but TTFT is dominated by prefill,
so speculative decoding helps less here).  FP8 on H100 provides $\sim$2$\times$
prefill speedup over FP16.

\textbf{Capacity planning:}  Apply Little's law: each request spends ${\sim}100$\,ms in
prefill and ${\sim}3$\,s in decode (200 tokens at ${\sim}15$\,ms/token at high batch),
so at 10K QPS roughly $10{,}000 \times 3 \approx 30{,}000$ requests are in flight at any
moment.  Per-request KV-cache for the assumed GQA-8 7B model at ${\sim}1.2$K tokens with
INT4 KV is $2 \cdot L \cdot H_{\text{kv}} \cdot S \cdot d_h \cdot b = 2 \times 32 \times
8 \times 1200 \times 128 \times 0.5 \approx 40$\,MB (the same cache with MHA in FP16
would be ${\sim}630$\,MB).  Cluster-wide that is ${\sim}1.2$\,TB of KV---only a few GB
per GPU spread across ${\sim}300$ GPUs.  With PagedAttention this is manageable; in this
regime the binding constraint is prefill compute, not KV memory.

\redflags
\begin{itemize}
    \item Jumping to ``just add more GPUs'' without reasoning about per-GPU capacity.
    \item Not mentioning continuous batching or PagedAttention.
    \item Ignoring the prefill vs.\ decode distinction for TTFT.
    \item Not quantifying anything---giving a purely qualitative answer.
\end{itemize}

\interviewtest{Whether the candidate can design a complete serving system with concrete
numbers, distinguishing prefill-bound TTFT from decode-bound token generation, and
reasoning about GPU-level capacity.}

\goodgreat
\begin{itemize}
    \item \badgeLfive{L5}: Mentions continuous batching, quantization, and prefix
          caching.  Gives rough GPU count estimates.
    \item \badgeLsix{L6}: Quantifies KV-cache memory, prefill latency, and per-GPU
          throughput.  Distinguishes prefill vs.\ decode bottlenecks.  Discusses
          failover and load balancing.
    \item \badgeLseven{L7}: Designs the full system including auto-scaling, model
          routing (small model for easy queries), graceful degradation under load,
          multi-region deployment, and cost analysis per query.
\end{itemize}

\followups{What changes if the average prompt is 32K tokens instead of 1K?  How would
you handle a 10$\times$ traffic spike?  What is the cost per query?}
\end{interviewq}

\vspace{0.5em}

% --- Question 2: KV-Cache ---
\begin{interviewq}
\textbf{Q: Explain the KV-cache.  Why is it necessary for autoregressive generation?
How much memory does it use?}

\badgeLfive{L5} \quad \badgetype{Conceptual} \quad \badgefreq{\freqhigh}

\stronganswer
In autoregressive generation, each new token attends to all previous tokens.  Without
caching, generating token $t$ requires recomputing the key and value projections for
all $t-1$ previous tokens---redundant work that grows quadratically over the full
sequence.

The KV-cache stores the key and value tensors from all previous positions so they are
computed only once.  Generation has two phases: \textbf{prefill} (process the entire
prompt in parallel, building the initial KV-cache---compute-bound) and \textbf{decode}
(generate one token at a time, appending to the cache---memory-bandwidth-bound).

\textbf{Memory calculation:}
$M_{\text{KV}} = 2 \times L \times H_{\text{kv}} \times S \times d_h \times b$.
For LLaMA-2 7B ($L{=}32$, $H_{\text{kv}}{=}32$, $d_h{=}128$) at 2K context in FP16:
$2 \times 32 \times 32 \times 2048 \times 128 \times 2 = 1.07$\,GB per request.  With
GQA ($H_{\text{kv}}{=}8$): 268\,MB.  This is why GQA is critical for production---it
reduces KV-cache memory by $4\times$ with minimal quality loss.

The KV-cache is the main memory bottleneck at inference time.  For long contexts (128K+
tokens), it can exceed the model weights in memory.  This is why PagedAttention (virtual
memory for KV-cache) and KV-cache quantization are active areas of optimization.

\redflags
\begin{itemize}
    \item Not knowing what the KV-cache stores or why it exists.
    \item Confusing prefill and decode phases.
    \item Inability to estimate memory usage with concrete numbers.
    \item Thinking KV-cache is optional or a ``nice to have.''
\end{itemize}

\interviewtest{Whether the candidate understands the fundamental inference mechanics of
autoregressive transformers and can reason quantitatively about memory.}

\goodgreat
\begin{itemize}
    \item \badgeLfive{L5}: Explains the two-phase structure, gives the memory formula,
          and computes a concrete example.
    \item \badgeLsix{L6}: Discusses GQA/MQA impact on KV-cache size, PagedAttention,
          prefix caching, and the distinction between compute-bound prefill and
          memory-bandwidth-bound decode.
    \item \badgeLseven{L7}: Discusses KV-cache quantization (INT8/INT4 for K/V tensors),
          KV-cache compression via attention sinks, and the interaction between
          KV-cache size and continuous batching capacity.
\end{itemize}

\followups{What is PagedAttention?  How does GQA reduce KV-cache size?  What happens to
KV-cache memory at 128K context?  How does continuous batching interact with KV-cache
management?}
\end{interviewq}

\vspace{0.5em}

% --- Question 3: Speculative Decoding ---
\begin{interviewq}
\textbf{Q: What is speculative decoding?  Why does it produce the same distribution as
standard autoregressive decoding?}

\badgeLsix{L6} \quad \badgetype{Conceptual} \quad \badgefreq{\freqmed}

\stronganswer
Speculative decoding accelerates autoregressive generation by using a small, fast
\emph{draft model} to propose candidate tokens that a large \emph{target model} verifies
in parallel.

The algorithm: (1) the draft model generates $k$ candidate tokens autoregressively (cheap:
$\sim$2ms per token for a 1B model).  (2) The target model processes all $k$ candidates
in a single forward pass (parallel, like prefill).  (3) At each position, compare the
draft model's token probability $q(x)$ to the target model's probability $p(x)$.  If
$q(x) \leq p(x)$, accept.  If $q(x) > p(x)$, accept with probability $p(x)/q(x)$,
reject otherwise.  On rejection, sample a correction from the residual distribution
$\max(0, p(x) - q(x))$ (normalized).

\textbf{Why the distribution is preserved:}  The accept/reject scheme is a form of
rejection sampling.  For any token $x$, the probability of it being accepted and output
is exactly $p(x)$:  If $q(x) \leq p(x)$, the token is always accepted (probability
$q(x)$ of being drafted $\times$ probability 1 of acceptance).  If $q(x) > p(x)$, it is
accepted with probability $p(x)/q(x)$ (so $q(x) \cdot p(x)/q(x) = p(x)$).  The
correction sampling handles the rejection case.  The net result is that every output
token is sampled exactly from $p(\cdot)$.

The speedup comes from the asymmetry between generation and verification: generating
$k$ tokens costs $k$ sequential forward passes, but verifying $k$ tokens costs only
one parallel forward pass.

\redflags
\begin{itemize}
    \item Claiming speculative decoding introduces quality loss.
    \item Not explaining the accept/reject mechanism.
    \item Confusing speculative decoding with beam search or other search methods.
    \item Not understanding why verification is cheaper than generation.
\end{itemize}

\interviewtest{Understanding of autoregressive generation mechanics, and whether the
candidate can reason about the mathematical guarantee of distribution preservation.}

\goodgreat
\begin{itemize}
    \item \badgeLfive{L5}: Correctly describes the draft-verify-accept/reject pipeline
          and states that the output distribution is preserved.
    \item \badgeLsix{L6}: Explains the rejection sampling argument formally, quantifies
          acceptance rates and speedup, discusses draft model selection.
    \item \badgeLseven{L7}: Discusses variants (Medusa, EAGLE, self-speculative), the
          theoretical speedup bound as a function of acceptance rate, and practical
          deployment considerations (batch interaction, variable acceptance rates).
\end{itemize}

\followups{How do you choose the draft model?  What is the theoretical maximum speedup?
How does speculative decoding interact with batching?  What is Medusa?}
\end{interviewq}

\vspace{0.5em}

% --- Question 4: Quantization Comparison ---
\begin{interviewq}
\textbf{Q: Compare GPTQ, AWQ, and INT8 quantization.  When would you use each?}

\badgeLsix{L6} \quad \badgetype{Trade-off} \quad \badgefreq{\freqmed}

\stronganswer
All three are post-training quantization methods, but they differ in precision,
mechanism, and quality-speed tradeoff.

\textbf{INT8 (e.g., LLM.int8(), SmoothQuant):}  Quantizes weights to 8-bit integers.
SmoothQuant migrates quantization difficulty from activations to weights by mathematically
equivalent channel-wise scaling.  $\sim$$2\times$ memory reduction, $\sim$$2\times$ speedup,
$<$1\% quality loss on nearly all benchmarks.  \emph{Use when:} you need a safe,
reliable compression with minimal risk---the conservative choice.

\textbf{GPTQ (4-bit):}  Layer-wise quantization using second-order (Hessian) information.
Quantizes each weight column while compensating for the error in remaining columns.
$\sim$$4\times$ memory reduction, $\sim$$3\times$ speedup, 1--3\% quality loss.  Slower
to quantize (hours for large models) but high-quality results.  \emph{Use when:} you need
4-bit compression and have time for careful quantization; good for offline preparation.

\textbf{AWQ (4-bit):}  Scales important weight channels (identified by activation
magnitude) before quantization.  Simpler and faster to apply than GPTQ, often matches or
exceeds GPTQ quality.  Same memory/speed profile.  \emph{Use when:} you need 4-bit
compression with fast quantization; increasingly the preferred 4-bit method.

\textbf{Decision framework:}
\begin{enumerate}
    \item If the model fits in memory at INT8 and quality is critical $\to$ INT8.
    \item If the model does not fit at INT8 or you need maximum throughput $\to$ AWQ
          (default 4-bit choice) or GPTQ.
    \item If you are on H100+ hardware $\to$ consider FP8 as a drop-in replacement for
          INT8 with potentially better quality.
\end{enumerate}

\redflags
\begin{itemize}
    \item Treating all quantization methods as interchangeable.
    \item Not knowing the mechanism behind GPTQ or AWQ (just ``it makes things smaller'').
    \item Not mentioning that quality loss varies by task and must be evaluated.
\end{itemize}

\interviewtest{Whether the candidate understands quantization at a mechanistic level
and can make principled deployment decisions.}

\goodgreat
\begin{itemize}
    \item \badgeLfive{L5}: Correctly describes all three methods and their tradeoffs.
    \item \badgeLsix{L6}: Explains the mechanisms (Hessian compensation for GPTQ,
          activation awareness for AWQ), discusses per-task evaluation necessity, and
          mentions FP8 as an emerging option.
    \item \badgeLseven{L7}: Discusses quantization of KV-cache separately from weights,
          mixed-precision strategies (4-bit weights + 8-bit KV-cache), and the
          interaction between quantization and other optimizations (speculative decoding,
          tensor parallelism).
\end{itemize}

\followups{How does GPTQ use the Hessian?  Why does AWQ focus on activation magnitude?
Can you quantize the KV-cache separately?  What is FP8 and when would you use it?}
\end{interviewq}

\vspace{0.5em}

% --- Question 5: Latency Reduction ---
\begin{interviewq}
\textbf{Q: How would you reduce the latency of an NLP inference system by 5$\times$?}

\badgeLfive{L5} \quad \badgetype{System Design} \quad \badgefreq{\freqhigh}

\stronganswer
A $5\times$ reduction requires stacking multiple optimizations.  I would approach this
systematically, starting with the highest-impact changes.

\textbf{1. Quantization ($2$--$4\times$):}  Move from FP16 to INT4 (AWQ/GPTQ).  This
alone can provide $\sim$$3\times$ speedup for memory-bandwidth-bound decoding, as the
model weights are $4\times$ smaller and read $4\times$ faster.

\textbf{2. Distillation or model selection ($2$--$6\times$):}  If the task allows it,
switch from a 70B model to a fine-tuned 7B model.  The smaller model is $\sim$$10\times$
faster, and with task-specific fine-tuning, quality loss can be minimal for narrow tasks.

\textbf{3. Speculative decoding ($2$--$3\times$):}  Reduces inter-token latency (but not
TTFT) by generating multiple tokens per target-model forward pass.

\textbf{4. Model compilation ($1.3$--$2\times$):}  TensorRT-LLM or \texttt{torch.compile}
fuse operations and optimize kernel selection.

\textbf{5. Serving infrastructure:}  Flash Attention (reduces prefill time), prefix
caching (eliminates redundant prefill), continuous batching (improves throughput and
reduces queuing delay).

\textbf{6. Architecture changes:}  Sliding window attention for long contexts, GQA to
reduce KV-cache bandwidth.

In practice, combining quantization ($3\times$) + compilation ($1.5\times$) already
yields $\sim$$4.5\times$.  Adding speculative decoding for decode-heavy workloads pushes
past $5\times$.  If the quality budget allows a smaller model, the gain is even larger.

\redflags
\begin{itemize}
    \item Giving only one technique (e.g., ``just quantize it'').
    \item Not stacking optimizations with multiplicative reasoning.
    \item Not distinguishing between TTFT and ITL optimization.
    \item Suggesting optimizations without estimating their individual contribution.
\end{itemize}

\interviewtest{Systematic optimization thinking---can the candidate identify multiple
independent axes of improvement and reason about their combined effect?}

\goodgreat
\begin{itemize}
    \item \badgeLfive{L5}: Lists 3+ techniques with approximate speedups and explains
          why they stack multiplicatively.
    \item \badgeLsix{L6}: Distinguishes TTFT vs.\ ITL optimizations, discusses
          quality-latency Pareto frontier, includes model selection as an option.
    \item \badgeLseven{L7}: Designs a full optimization roadmap with measurement at each
          stage, discusses diminishing returns, considers cost vs.\ latency tradeoff,
          and mentions model routing as a system-level optimization.
\end{itemize}

\followups{Which of these optimizations affect TTFT vs.\ decode latency?  How would you
measure the quality impact of each?  What if you cannot change the model?}
\end{interviewq}

\vspace{0.5em}

% --- Question 6: BPE Tokenization ---
\begin{interviewq}
\textbf{Q: Explain BPE tokenization.  How is it trained?  How does it handle OOV words?}

\badgeLfive{L5} \quad \badgetype{Conceptual} \quad \badgefreq{\freqhigh}

\stronganswer
Byte Pair Encoding is a subword tokenization algorithm that learns a vocabulary of
variable-length tokens from a training corpus.

\textbf{Training:}  (1) Initialize the vocabulary with all individual characters (or
bytes, in byte-level BPE).  (2) Count the frequency of every adjacent token pair in the
corpus.  (3) Merge the most frequent pair into a new token and add it to the vocabulary.
(4) Repeat steps 2--3 for a predefined number of merges (e.g., 32,000 merges for a 32K
vocabulary).  The ordered list of merge rules \emph{is} the tokenizer.

\textbf{Encoding (inference):}  Given a new string, start with individual characters and
apply the learned merge rules greedily in priority order.  The result is a sequence of
subword tokens.

\textbf{OOV handling:}  BPE \emph{never} produces an unknown token.  Any unseen word is
decomposed into known subword units.  In byte-level BPE (used by GPT-2+), the base
vocabulary is the 256 possible byte values, so \emph{any} byte sequence has a valid
tokenization.  A rare word like ``unparalleledly'' might become
\texttt{["un", "parallel", "ed", "ly"]}---the model retains morphological information
even for words it has never seen as a complete unit.

\textbf{Vocabulary size tradeoff:}  Larger vocabulary $\to$ shorter sequences (less
compute) but larger embedding tables.  32K--100K is typical; modern models trend toward
100K+ because the attention savings from shorter sequences outweigh the embedding cost.

\redflags
\begin{itemize}
    \item Not knowing the training algorithm (just ``it splits words into pieces'').
    \item Saying BPE produces \texttt{[UNK]} tokens.
    \item Confusing BPE with WordPiece or not knowing the difference.
    \item Not mentioning the vocabulary size tradeoff.
\end{itemize}

\interviewtest{Whether the candidate understands tokenization as an algorithm, not just
a preprocessing step---and whether they appreciate its impact on model behavior.}

\goodgreat
\begin{itemize}
    \item \badgeLfive{L5}: Describes the full training and encoding algorithm, explains
          OOV handling, and mentions the vocabulary size tradeoff.
    \item \badgeLsix{L6}: Compares BPE with WordPiece and Unigram, discusses byte-level
          BPE for multilingual handling, and explains why LLMs struggle with
          character-level tasks (tokenization breaks character boundaries).
    \item \badgeLseven{L7}: Discusses the tokenizer tax for non-English languages,
          tokenizer training as a separate optimization problem from model training,
          and the impact of tokenizer design on cross-lingual transfer.
\end{itemize}

\followups{What is the difference between BPE and WordPiece?  Why do LLMs struggle
with counting letters in words?  How does tokenizer choice affect multilingual
performance?}
\end{interviewq}

\vspace{0.5em}

% --- Question 7: Distillation vs Quantization vs Pruning ---
\begin{interviewq}
\textbf{Q: When would you choose knowledge distillation over quantization or pruning to
compress a model?}

\badgeLfive{L5} \quad \badgetype{Trade-off} \quad \badgefreq{\freqmed}

\stronganswer
These three compression techniques are complementary, not competing.  Each has a
different profile:

\textbf{Quantization:}  Reduces numerical precision.  $2$--$4\times$ smaller, minimal
quality loss, no training required (PTQ).  Choose quantization as the \emph{default first
step}---it is the lowest-effort, highest-reward compression.

\textbf{Pruning:}  Removes weights or structures.  Structured pruning (removing heads or
layers) gives real speedup without hardware constraints; unstructured pruning requires
sparse hardware support.  Choose pruning when you need a smaller \emph{dense} model
(structured) or when targeting hardware with sparse acceleration (unstructured).

\textbf{Distillation:}  Trains a new, smaller architecture.  Highest potential
compression ($5$--$10\times$ with architecture change) but requires a full training run.
Choose distillation when: (1) you need $>$$4\times$ compression (beyond what quantization
provides), (2) you want to change the architecture (e.g., 12-layer $\to$ 6-layer), or
(3) you have a strong teacher model and sufficient compute for student training.

\textbf{The practical order:}
\begin{enumerate}
    \item Start with \textbf{quantization} (INT8 or INT4)---immediate gain, zero training.
    \item If more compression is needed, \textbf{prune} unnecessary heads/layers.
    \item If quality is still insufficient at the target size, \textbf{distill} into a
          properly sized student.
    \item Apply quantization to the distilled/pruned model for additional compression.
\end{enumerate}

\redflags
\begin{itemize}
    \item Treating the three as mutually exclusive.
    \item Recommending distillation when quantization would suffice (wasteful).
    \item Not mentioning that distillation requires training compute.
\end{itemize}

\interviewtest{Practical engineering judgment about when to use which compression
technique, and awareness that they compose.}

\goodgreat
\begin{itemize}
    \item \badgeLfive{L5}: Correctly describes all three techniques, gives scenarios
          for each, and notes they can be combined.
    \item \badgeLsix{L6}: Discusses the practical order of application, mentions
          specific methods (GPTQ, DistilBERT), and quantifies typical compression
          ratios for each technique.
    \item \badgeLseven{L7}: Discusses the quality-efficiency Pareto frontier, mentions
          that distillation can use synthetic data from the teacher (no labeled data
          needed), and considers the full lifecycle cost (training cost of distillation
          vs.\ ongoing inference cost savings).
\end{itemize}

\followups{Can you quantize a distilled model?  How does structured pruning differ from
distillation?  What is the cost-benefit analysis of distillation (training cost vs.\
inference savings)?}
\end{interviewq}

\vspace{0.5em}

% --- Question 8: Tokenizer Tax ---
\begin{interviewq}
\textbf{Q: How do you handle the ``tokenizer tax'' when deploying NLP systems for
multilingual users?}

\badgeLsix{L6} \quad \badgetype{System Design} \quad \badgefreq{\freqmed}

\stronganswer
The tokenizer tax refers to the phenomenon where non-English text (especially non-Latin
scripts) is tokenized into significantly more tokens than equivalent English text.  For
example, the same sentence might be 7 tokens in English but 30 tokens in Hindi with a
tokenizer trained primarily on English.  This has four cascading effects: shorter
effective context windows, higher compute cost, higher API cost, and slower generation.

\textbf{Mitigation at the tokenizer level:}
\begin{itemize}
    \item Train the tokenizer on a \emph{balanced} multilingual corpus, ensuring
          non-English languages have adequate representation.  LLaMA 3's tokenizer was
          explicitly redesigned for better multilingual fertility.
    \item Increase vocabulary size (100K--150K tokens) to include more non-English
          subwords.  This reduces fertility for non-English languages at the cost of
          a larger embedding table.
    \item Use byte-level BPE (ensures universal coverage) combined with multilingual
          training data.
    \item Measure \textbf{fertility} (tokens per word) across all target languages as a
          tokenizer evaluation metric.  Flag any language with fertility $>$$2\times$
          the English baseline.
\end{itemize}

\textbf{Mitigation at the system level:}
\begin{itemize}
    \item For API-based systems, consider per-character or per-word pricing rather than
          per-token to avoid penalizing non-English users.
    \item Allocate larger context budgets for languages with higher fertility.
    \item Use language-specific adapter models or language-specific fine-tuning to
          compensate for reduced per-token quality in underrepresented languages.
    \item Monitor quality metrics \emph{per language}, not just on average.
\end{itemize}

\redflags
\begin{itemize}
    \item Not knowing what the tokenizer tax is.
    \item Suggesting ``just use a bigger context window'' without addressing the root cause.
    \item Ignoring the cost and latency implications.
    \item Not mentioning tokenizer training as the primary lever.
\end{itemize}

\interviewtest{Awareness of multilingual deployment challenges and ability to address
them at both the model and system level.}

\goodgreat
\begin{itemize}
    \item \badgeLfive{L5}: Explains the tokenizer tax with concrete examples and
          suggests training on balanced multilingual data.
    \item \badgeLsix{L6}: Discusses vocabulary size tradeoffs, fertility as a metric,
          and system-level mitigations (pricing, context budgets, per-language monitoring).
    \item \badgeLseven{L7}: Discusses cross-lingual transfer failure modes (script
          mismatch, morphological mismatch), the tension between shared vocabulary
          and per-language quality, and strategies like language-adaptive tokenization
          or language mixture training schedules.
\end{itemize}

\followups{How would you evaluate tokenizer quality across 50 languages?  What is the
relationship between tokenizer fertility and downstream quality?  How does tokenizer
design affect cross-lingual transfer?}
\end{interviewq}

\vspace{0.5em}

% --- Question 9: Production Monitoring ---
\begin{interviewq}
\textbf{Q: You have deployed an LLM-powered feature and users report quality has degraded
over the past week, but your automated metrics show no change.  How do you diagnose this?}

\badgeLsix{L6} \quad \badgetype{System Design} \quad \badgefreq{\freqmed}

\stronganswer
This is a classic case where automated metrics diverge from user experience.  I would
investigate along several axes.

\textbf{1. Input distribution shift:}  Compare the input distribution from the past week
to the baseline period.  Check average prompt length, topic distribution (via clustering
input embeddings), language mix, and any new input patterns.  Users may be asking
different types of questions that the model handles poorly, even though the model itself
has not changed.

\textbf{2. Metric blind spots:}  Automated metrics may not capture what users care about.
If we are monitoring perplexity or BLEU, these correlate poorly with user-perceived quality
for open-ended generation.  Check: are we measuring factuality?  Helpfulness?  Tone?
Add LLM-as-judge evaluation on a sample of recent traffic.

\textbf{3. System-level issues:}  Check for silent failures: increased timeout rates
(responses truncated), load balancer changes routing to a different model version,
cache corruption returning stale responses, or infrastructure changes affecting context
length limits.

\textbf{4. Upstream changes:}  Check for changes in the retrieval system (if RAG),
system prompt, safety filters, or post-processing logic.  Any of these can degrade
output quality without the LLM itself changing.

\textbf{5. User expectations:}  Consider whether user expectations have shifted (e.g.,
competing products improved, creating a relative quality perception change).

\textbf{Diagnosis process:}  Sample 50--100 recent user complaints.  Manually evaluate
them alongside model outputs from both periods.  Categorize failure modes (factual errors,
formatting issues, refusals, irrelevant responses).  This often immediately reveals the
root cause.

\redflags
\begin{itemize}
    \item Assuming the model must have changed if quality degraded.
    \item Not considering input distribution shift.
    \item Relying solely on automated metrics without manual inspection.
    \item Not considering system-level or upstream changes.
\end{itemize}

\interviewtest{Production debugging mindset---can the candidate systematically reason
about failure modes in a complex system beyond just the model?}

\goodgreat
\begin{itemize}
    \item \badgeLfive{L5}: Checks input distribution, model version, and automated
          metric coverage.  Samples complaints manually.
    \item \badgeLsix{L6}: Considers the full system (retrieval, prompts, filters, infra),
          proposes specific monitoring additions, and designs a structured diagnosis
          workflow.
    \item \badgeLseven{L7}: Designs a comprehensive observability stack with
          embedding-space drift detection, LLM-as-judge pipelines, user feedback
          integration, and automatic alerting on quality segments (per language,
          per topic, per user cohort).
\end{itemize}

\followups{How would you set up monitoring to catch this proactively?  How do you
measure ``quality'' for open-ended generation?  What is the role of human evaluation
in production systems?}
\end{interviewq}

\vspace{0.5em}

% --- Question 10: Model Routing ---
\begin{interviewq}
\textbf{Q: Describe a model routing strategy that reduces serving cost by 5$\times$ while
maintaining quality.}

\badgeLsix{L6} \quad \badgetype{System Design} \quad \badgefreq{\freqlow}

\stronganswer
The key insight is that most production queries are ``easy'' and do not require a
frontier model.  A model routing system directs each query to the cheapest model capable
of handling it well.

\textbf{Architecture:}
\begin{enumerate}
    \item A lightweight \textbf{router} (small classifier or embedding model) evaluates
          each incoming query.
    \item \textbf{Easy queries} (simple factual QA, routine classification, formatting
          requests) route to a small model (e.g., 7B INT4)---$\sim$$50\times$ cheaper
          than GPT-4.
    \item \textbf{Hard queries} (complex reasoning, multi-step tasks, ambiguous requests)
          route to a large model (e.g., 70B FP16 or API call to frontier model).
    \item If 80\% of traffic is routable to the small model (typical for many workloads),
          effective cost is: $0.8 \times c_{\text{small}} + 0.2 \times c_{\text{large}}$.
          With $c_{\text{small}} = c_{\text{large}} / 50$, this is $0.016c + 0.2c =
          0.216c$---a $\sim$$4.6\times$ cost reduction.
\end{enumerate}

\textbf{Router design:}
\begin{itemize}
    \item Train a difficulty classifier on a dataset of (query, difficulty) labels, where
          difficulty is determined by comparing small-model and large-model outputs.
    \item Features: query length, embedding, presence of reasoning keywords, topic.
    \item Alternatively: always try the small model first; if the response confidence is
          low or fails a quality check, fall back to the large model.
\end{itemize}

\textbf{Quality guarantee:}  Monitor per-route quality metrics.  If the small model's
quality degrades on any query segment, the router thresholds can be adjusted dynamically
to route more traffic to the large model.

\redflags
\begin{itemize}
    \item Proposing only a single model without considering routing.
    \item Not quantifying the cost savings.
    \item Not addressing how to maintain quality for routed queries.
\end{itemize}

\interviewtest{System-level cost optimization thinking and practical deployment
architecture.}

\goodgreat
\begin{itemize}
    \item \badgeLfive{L5}: Describes the basic routing concept and gives a cost
          calculation.
    \item \badgeLsix{L6}: Designs the router architecture, discusses fallback strategies,
          and proposes quality monitoring per route.
    \item \badgeLseven{L7}: Discusses cascading routers (multiple tiers), learned routing
          policies, the interaction between routing and batching efficiency, and A/B
          testing the router itself.
\end{itemize}

\followups{How would you train the difficulty classifier?  What if the small model fails
silently (confidently wrong)?  How does routing interact with latency SLAs?}
\end{interviewq}


% ============================================================================
% CHAPTER SUMMARY
% ============================================================================
\section*{Chapter Summary}

\begin{enumerate}
    \item \textbf{Tokenization} determines vocabulary size, sequence length, and
          multilingual performance.  BPE (bottom-up merging by frequency), WordPiece
          (merging by likelihood), and Unigram (top-down pruning by likelihood loss)
          are the three core algorithms.  Byte-level BPE eliminates UNK tokens entirely.

    \item \textbf{KV-cache} is the essential inference optimization for autoregressive
          transformers.  Memory scales as $2 \times L \times H_{\text{kv}} \times S
          \times d_h \times b$ and is the primary bottleneck for long-context serving.
          GQA reduces KV-cache memory by $4$--$8\times$.

    \item \textbf{Continuous batching} and \textbf{PagedAttention} are the infrastructure
          primitives that make production LLM serving viable, improving throughput by
          $2$--$5\times$ and memory utilization from $\sim$20\% to $\sim$95\%.

    \item \textbf{Quantization} (INT8, INT4 via GPTQ/AWQ) provides $2$--$4\times$
          memory and speed improvements with $<$3\% quality loss.  Always evaluate on
          your specific task.

    \item \textbf{Knowledge distillation} trains a small student to mimic a large
          teacher using soft labels.  The temperature-scaled KL-divergence loss
          transfers ``dark knowledge'' about inter-class relationships.

    \item \textbf{Speculative decoding} uses a small draft model to propose tokens
          verified by the target model in parallel, providing $2$--$3\times$ speedup
          with provably identical output distribution.

    \item \textbf{Multilingual deployment} requires awareness of the tokenizer tax,
          cross-lingual transfer limitations, and per-language quality monitoring.

    \item \textbf{MLOps} for NLP encompasses latency monitoring (TTFT, ITL, P99),
          quality tracking, data drift detection, cost optimization via model routing,
          and continuous feedback loops.
\end{enumerate}

\begin{keyinsight}[Chapter 12 Summary: What Interviewers Are Really Testing]
When interviewers ask production NLP questions, they are testing three things:
(1)~\emph{Do you think in numbers?}---Can you estimate KV-cache memory, prefill latency,
and throughput for a specific model?  Candidates who reason quantitatively signal that
they have actually deployed models, not just trained them.
(2)~\emph{Do you know the full stack?}---Tokenization, serving, optimization, and
monitoring are all part of production NLP.  Depth in any one of these is valuable, but
breadth across all of them signals a staff-level engineer.
(3)~\emph{Can you make tradeoffs?}---Every optimization trades something (quality, cost,
complexity, generality).  The best candidates do not just list techniques---they reason
about when each technique is appropriate and when it is not.
\end{keyinsight}
